% ============================================================
% Preamble for the master thesis proposal.
% Edit fonts, margins and macros here, not in proposal.tex.
% ============================================================

% ---------- encoding and fonts ----------
\usepackage[T1]{fontenc}
\usepackage[utf8]{inputenc}
\usepackage{lmodern}
\usepackage{microtype}

% ---------- page layout ----------
\usepackage[a4paper, margin=27mm]{geometry}
\usepackage{setspace}
\onehalfspacing
\usepackage{parskip}

% ---------- mathematics ----------
\usepackage{amsmath, amssymb, amsthm, mathtools}
\usepackage{bm}
\usepackage{siunitx}          % \SI{9.81}{\meter\per\second\squared}
\sisetup{separate-uncertainty = true}

% ---------- figures and tables ----------
\usepackage{graphicx}
\graphicspath{{figures/}}
\usepackage{subcaption}
\usepackage{booktabs}
\usepackage{array}
\usepackage{longtable}
\usepackage[font=small, labelfont=bf]{caption}

% ---------- lists ----------
\usepackage{enumitem}
\setlist{itemsep=3pt, topsep=5pt}

% ---------- code listings ----------
\usepackage{listings}
\lstset{
  basicstyle=\ttfamily\small,
  breaklines=true,
  frame=lines,               % thin rules above and below, no side lines
  rulecolor=\color{black!30},
  captionpos=b,
  showstringspaces=false
}

% ---------- colour (kept minimal) ----------
\usepackage[dvipsnames]{xcolor}

% ---------- headers: clean academic, no rules, just page numbers ----------
\pagestyle{plain}

% ---------- section styling ----------
\usepackage{titlesec}
\titleformat{\section}{\large\bfseries}{\thesection}{0.7em}{}
\titlespacing{\section}{0pt}{18pt}{7pt}
\titleformat{\subsection}{\bfseries}{\thesubsection}{0.6em}{}
\titlespacing{\subsection}{0pt}{13pt}{4pt}

% ---------- theorem-like environments ----------
\theoremstyle{plain}
\newtheorem{theorem}{Theorem}[section]
\newtheorem{proposition}[theorem]{Proposition}
\theoremstyle{definition}
\newtheorem{definition}[theorem]{Definition}
\newtheorem{example}[theorem]{Example}
\theoremstyle{remark}
\newtheorem{remark}[theorem]{Remark}

% ---------- boxes for highlighted statements ----------
\usepackage{mdframed}
\mdfsetup{skipabove=12pt, skipbelow=12pt}
% keybox: thin horizontal rules above and below only, no vertical sides.
\newenvironment{keybox}
  {\begin{mdframed}[linewidth=0.5pt, topline=true, bottomline=true,
     leftline=false, rightline=false,
     innertopmargin=8pt, innerbottommargin=10pt,
     innerleftmargin=0pt, innerrightmargin=0pt]}
  {\end{mdframed}}

% ---------- bibliography ----------
% natbib + bibtex, both bundled with tectonic, so the document
% builds anywhere with no extra tools installed.
% Cite with \cite{key} or \citet{key} for "Author (year)" style.
\usepackage[numbers, sort&compress]{natbib}
\bibliographystyle{unsrtnat}      % numbered in order of first citation

% ---------- links and cross-references ----------
\usepackage[hidelinks]{hyperref}
\usepackage[capitalise, noabbrev]{cleveref}

% ============================================================
% Drafting helpers. Set \draftmode to false before submitting:
% notes and the TODO list then disappear from the output.
% ============================================================
\newif\ifdraftmode
\draftmodetrue          % <-- set to \draftmodefalse for the final version

\usepackage{xparse}
\NewDocumentCommand{\todo}{m}{%
  \ifdraftmode\textcolor{BrickRed}{\textbf{[TODO: #1]}}\fi}
\NewDocumentCommand{\note}{m}{%
  \ifdraftmode\textcolor{RoyalBlue}{\textit{[note: #1]}}\fi}

% ============================================================
% Notation macros. Add your own here so the notation stays
% consistent across sections; change it once, it changes
% everywhere.
% ============================================================
\newcommand{\R}{\mathbb{R}}
\newcommand{\vect}[1]{\bm{#1}}              % vectors: \vect{u}
\newcommand{\mat}[1]{\bm{#1}}               % matrices: \mat{A}
\newcommand{\norm}[1]{\left\lVert #1 \right\rVert}
\newcommand{\abs}[1]{\left\lvert #1 \right\rvert}
\newcommand{\inner}[2]{\left\langle #1, #2 \right\rangle}
\newcommand{\T}{\mathsf{T}}                 % transpose: A^\T
\newcommand{\inv}{^{-1}}
\newcommand{\ket}[1]{\lvert #1 \rangle}     % quantum notation
\newcommand{\bra}[1]{\langle #1 \rvert}
\newcommand{\braket}[2]{\langle #1 \vert #2 \rangle}
\newcommand{\op}[1]{\hat{#1}}
\newcommand{\pd}[2]{\frac{\partial #1}{\partial #2}}
\newcommand{\dd}{\mathrm{d}}
\DeclareMathOperator{\diver}{div}
\DeclareMathOperator*{\argmin}{arg\,min}
\DeclareMathOperator{\rank}{rank}
\DeclareMathOperator{\diag}{diag}
